\section{Край магнита}
\label{sec:edge}

Поле реального диполя на торце спадает, и по уравнениям Максвелла при этом появляется продольная компонента $B_s$. В пределе жёсткого края она становится $\delta$-функцией и даёт тонкий вертикальный удар, пропорциональный горизонтальному наклону траектории. На дисперсионной орбите наклон равен $D_1'\delta$, т.\,е. частица с $\delta\neq0$ входит в магнит с перпендикулярным торцом под эффективным углом. Это обычная краевая фокусировка, но хроматическая. В главе выведен её вклад в $\xi_y$, объяснено, почему в $\xi_x$ его нет, и описан переход от конечного краевого поля к жёсткому краю.

\subsection{Продольное поле жёсткого края}

Пусть вход в диполь находится при $s=0$, торец перпендикулярен орбите, и в средней плоскости $B_y=B_0\Theta(s)$. В вакууме $\operatorname{rot}\mathbf{B}=0$, и $s$-компонента условия $\partial B_s/\partial y=\partial B_y/\partial s$ в первом порядке по $y$ даёт
\begin{equation}
\label{eq:Bs}
B_s=y\,B_0\,\delta_D(s)\quad\text{(вход)},
\qquad
B_s=-y\,B_0\,\delta_D(s-L)\quad\text{(выход)},
\end{equation}
где $\delta_D$ --- дельта-функция Дирака, $L$ --- длина магнита. Кривизной орбиты на бесконечно тонком слое можно пренебречь. Гамильтониан \eqref{eq:H} этой компоненты не содержит: он построен для поля, не зависящего от $s$.

\subsection{Краевой удар}

Сила Лоренца $\mathrm{d}\mathbf{P}/\mathrm{d}t=q\,\mathbf{v}\times\mathbf{B}$. Вертикальная компонента $(\mathbf{v}\times\mathbf{B})_y=v_sB_x-v_xB_s$; в слое $B_x$ не имеет особенности, остаётся член с $B_s$. Интегрирование по слою ($\mathrm{d}t=\mathrm{d}s/v_s$, $v_x/v_s=x'$) и деление на $p_s$ дают
\begin{equation}
\label{eq:kick}
\Delta p_y=-\frac{x'\,y}{\rho}\quad\text{(вход)},
\qquad
\Delta p_y=+\frac{x'\,y}{\rho}\quad\text{(выход)},
\end{equation}
с $qB_0/p_s=1/\rho$. Горизонтальная компонента $(\mathbf{v}\times\mathbf{B})_x=v_yB_s-v_sB_y$ даёт от $B_s$ удар $\propto y'y/\rho$: он квадратичен по вертикальным переменным и от $X$ не зависит.

Это та же формула, что для торца, повёрнутого на угол $e$: вертикальная краевая фокусировка $\Delta p_y=-(\tan e/\rho)\,y$. Для перпендикулярного торца роль $\tan e$ играет наклон траектории $x'$: на входе эффективный угол $e+x'$, на выходе $e-x'$ (в принятой здесь конвенции знаков, в которой $e>0$ фокусирует по $y$). На референсной орбите $x'=0$, и линейной линзы нет. Частица с $\delta\neq0$ идёт по дисперсионной орбите с $x'=D_1'\delta$ и видит тонкую вертикальную линзу.

\subsection{Вклад в хроматичность}

Подстановка $x'=D_1'\delta$ в \eqref{eq:kick} даёт на входе $\Delta p_y=-k\,y$ с $k=D_1'\delta/\rho$ (фокусировка), на выходе $k=-D_1'\delta/\rho$. Тонкая линза сдвигает число колебаний на $\Delta\nu_y=\beta_yk/(4\pi)$, и
\begin{equation}
\label{eq:xiF}
\xi^{\mathrm{F}}_y=\frac{1}{4\pi\rho}\sum_{\text{диполи}}
\Bigl[\beta_yD_1'\big|_{\text{вход}}-\beta_yD_1'\big|_{\text{выход}}\Bigr],
\qquad
\xi^{\mathrm{F}}_x=0.
\end{equation}
В гамильтоновом виде это член $\mathcal{H}_2^{(1)}=\pm(D_1'/2\rho)\,Y^2\,\delta_D(s-s_{\text{торца}})$, который входит в \eqref{eq:xi-ham} так же, как остальные (раздел~\ref{sec:def}). Если оптика симметрична относительно середины диполя, $D_1'$ на входе и выходе имеют разные знаки, и вклады торцов складываются, а не сокращаются. Знак суммы определяется оптикой; при $D_1'>0$ на входе и $D_1'<0$ на выходе вклад положителен, т.\,е. уменьшает модуль отрицательной натуральной $\xi_y$.

\paragraph{Почему нет горизонтального вклада.}
Обычная горизонтальная краевая фокусировка $+(\tan e/\rho)\,x$ возникает потому, что при повёрнутом торце длина пути в поле зависит от $x$. При перпендикулярном торце граница поля --- это поверхность $s=\mathrm{const}$ сопровождающей системы, и путь в поле от $x$ не зависит; наклонная траектория описывается гамильтонианом тела \eqref{eq:H} точно, а её вклад уже содержится в членах C и E. Горизонтальная часть удара от $B_s$ \eqref{eq:kick} квадратична по $y$ и на $\xi_x$ в первом порядке не влияет.

\paragraph{Если торцы повёрнуты.}
При $e\neq0$ к \eqref{eq:xiF} добавляются хроматические части обычной краевой фокусировки: удары $\mp(\tan e/\rho)\,x$, $\pm(\tan e/\rho)\,y$ в канонических импульсах от $\delta$ не зависят, но к ним прибавляются члены от сдвига точки линеаризации (квадратичные коэффициенты края, умноженные на $D_1$) и от зависимости эффективного угла от $D_1'\delta$. Для магнитов этой заметки $e=0$, и подробно этот случай не разбирается.

\subsection{Конечное краевое поле и предел жёсткого края}

Реальный спад описывают функцией Энге~\cite{Enge1967,ValetovPhD}
\begin{equation}
\label{eq:enge}
B_y(s)=B_0F(s),
\qquad
F(s)=\frac{1}{1+\exp\bigl(a_1+a_2(s/g)+\dots+a_n(s/g)^n\bigr)},
\end{equation}
где $g$ --- апертура (зазор), задающая масштаб спада. Положение эквивалентного жёсткого края --- эффективная граница поля (EFB) --- выбирается из равенства площадей:
\begin{equation}
\label{eq:efb}
\int_{-\infty}^{\infty}\bigl[F(s)-\Theta(s_{\mathrm{EFB}}-s)\bigr]\,\mathrm{d}s=0
\end{equation}
(для выхода; для входа зеркально). Тогда интеграл поля по краю совпадает с жёстким краем в $s_{\mathrm{EFB}}$, и угол поворота сохраняется. В~\cite{ValetovPhD} то же условие записано через точки $z_{\mathrm{int}}$, $z_{\mathrm{ext}}$ внутри и вне элемента.

В том же порядке по $y$ продольное поле $B_s=yB_0F'(s)$, и вертикальный удар края
\begin{equation}
\label{eq:kick-fringe}
\Delta p_y=-\frac{1}{\rho}\int x'(s)\,y(s)\,F'(s)\,\mathrm{d}s .
\end{equation}
Если на длине спада $\sim g$ наклон $x'$ и координата $y$ почти постоянны, $\int F'\,\mathrm{d}s=\pm1$ и \eqref{eq:kick-fringe} переходит в \eqref{eq:kick}, а вклад в хроматичность --- в \eqref{eq:xiF} с $\beta_y$, $D_1'$ на EFB. Поправки к пределу жёсткого края имеют порядок $g$:
\begin{itemize}
\item наклон $x'$ и $y$ меняются внутри слоя (кривизна $F/\rho$ действует уже в краевом поле);
\item следующий порядок по $y$ даёт известную поправку к вертикальной краевой фокусировке $\propto g\,K_{\mathrm{fr}}/\rho^2$ (при $e=0$ --- дефокусирующую), где $K_{\mathrm{fr}}=\int F(1-F)\,\mathrm{d}s/g$ (параметр \texttt{FINT} в MAD-X); она сдвигает $\nu_y$ уже при $\delta=0$, а через $1/(1+\delta)$ даёт и малую хроматическую поправку.
\end{itemize}
При $g\to0$ все они исчезают линейно по $g$, и модель с конечным краем стремится к жёсткому краю, но не к модели <<только тело>>: удар \eqref{eq:kick} от формы спада не зависит, пока спад короткий. В~\cite{ValetovChrom2026} это видно по иерархии моделей кольца Muon g-2: жёсткий край на EFB и реалистичный спад Энге дают практически одинаковую хроматичность, а неучёт смещения EFB заметно её меняет.

Отметим, что связь бетатронного удлинения орбиты с хроматичностью $\Lambda=-(2\pi/C_s)(J_x\xi_x+J_y\xi_y)$~\cite{EqNote} от формы края не зависит: обе её части берутся из одного гамильтониана, и член \eqref{eq:xiF} входит в $\Lambda$ через полное $\xi_y$.
